\subsection{D. The letters and the printed EGA statements}

We distinguish a false assertion, an invalid proof of a true assertion,
an unnecessary hypothesis, and a convention.  In particular the letter
of 29 October concerns \emph{EGA IV} 19.5.5 and 19.5.6, not
\emph{EGA $0_{\rm IV}$}; the latter numbers denote different text.
The primary text of $0_{\rm IV}$ used below is the 1964 printing,
pp.~85--86, 94, and 124--127, opened both as a scan and as searchable
text.  The letters are cote 23, scans 044--065.

\subsubsection{D.1. A repair of $0_{\rm IV}$ 19.3.12}

\begin{lemma}[Proved: the linear compactness needed here]\label{SA23-D-1}
Let $R$ be a ring and let $(M_i)_{i\in I}$ be Artinian $R$-modules.
Consider any set of affine $R$-linear equations in variables
$z_i\in M_i$, each equation involving only finitely many variables
and taking values in an $R$-module.  If every finite subsystem has a
solution, then the entire system has a solution.
\end{lemma}
\begin{proof}
A family of cosets of submodules of one Artinian module, with the
finite intersection property, has nonempty intersection.  Indeed,
among the submodules underlying its nonempty finite intersections
choose a minimal one, using the descending chain condition.  Its
coset is unchanged on intersection with any further member of the
family, since two nonempty cosets of the same submodule, one contained
in the other, are equal.

Well-order the variables.  We choose their values successively while
preserving finite solvability of the original equations together with
the equations fixing all values already chosen.  For the next
variable $z_i$, every finite subsystem projects to a nonempty coset
of a submodule of $M_i$: the solution set of a finite affine linear
system is a coset, and so is its image under a linear projection.
These cosets have the finite intersection property by the induction
hypothesis.  Choose a value in their intersection by the preceding
paragraph.  At a limit ordinal finite solvability persists, since a
finite subsystem mentions only finitely many previous choices.  At
the end all the equations hold.  This proof does not assert any
analogous compactness for polynomial equations.
\end{proof}

\begin{lemma}[Proved: closed quotients of a pro-Artinian ring]\label{SA23-D-2}
Let $C$ be a topological ring with a directed fundamental system of
open ideals $(U_\lambda)$ such that $C_\lambda=C/U_\lambda$ is
Artinian and $C\simeq\varprojlim_\lambda C_\lambda$ topologically.
For every closed ideal $D$ of $C$, the canonical map
\[
 C/D\longrightarrow\varprojlim_\lambda C/(D+U_\lambda)
\]
is an isomorphism of topological rings.  In particular the closed
quotient is again complete and pro-Artinian.
\end{lemma}
\begin{proof}
Injectivity is the identity
$D=\bigcap_\lambda(D+U_\lambda)$, which is the definition of
closedness for this topology.  Given compatible residues
$d_\lambda\in C/(D+U_\lambda)$, use variables
$c_\lambda\in C_\lambda$ with the linear equations prescribing
their images to be $d_\lambda$ and requiring compatibility under the
quotient maps.  Each $C_\lambda$ is an Artinian $C$-module.  A finite
subsystem is solved by choosing an index dominating all its indices
and lifting the prescribed residue there to an element of $C$.
Lemma~\ref{SA23-D-1} gives compatible $c_\lambda$, hence an element
of $C$.  This proves surjectivity.  Both topologies have as a
fundamental system of open ideals the kernels of the displayed maps
to $C/(D+U_\lambda)$, proving the topological assertion.
\end{proof}

\begin{theorem}[Proved: the statement of $0_{\rm IV}$ 19.3.12 is true]\label{SA23-D-3}
Let $A$ be a topological ring, $B$ a formally smooth topological
$A$-algebra, and $C$ a topological $A$-algebra.  Suppose that
$C=\varprojlim C/U_\lambda$ topologically with each $C/U_\lambda$
Artinian.  Let $J\subset C$ be a closed ideal whose image in every
$C/U_\lambda$ is nilpotent; this follows in particular when $J$ is
topologically nilpotent as in the printed proposition.  Suppose also
that the square of every open ideal of $B$ is open.  Every continuous
$A$-algebra homomorphism $u:B\to C/J$ lifts to a continuous
$A$-algebra homomorphism $v:B\to C$.
\end{theorem}
\begin{proof}
We first prove the assertion when $J^2=0$.  Write
$J_\lambda=(J+U_\lambda)/U_\lambda$ and
$u_\lambda:B\to C_\lambda/J_\lambda$ for the induced map.
Formal smoothness gives an $A$-algebra lift
$w^0_\lambda:B\to C_\lambda$ for each $\lambda$.
The possible algebra lifts have the form
$w_\lambda=w^0_\lambda+z_\lambda$, where the variables
$z_\lambda(b)\in J_\lambda$ for $b\in B$ satisfy the equations
\[
\begin{split}
 z_\lambda(b+b')&=z_\lambda(b)+z_\lambda(b'),\\
 z_\lambda(bb')&=w^0_\lambda(b)z_\lambda(b')
                  +w^0_\lambda(b')z_\lambda(b),\\
 z_\lambda(a)&=0\quad(a\in A),\qquad z_\lambda(1)=0.
\end{split}
\]
Here $a$ means its image in $B$.  These are linear equations over
$C$ in Artinian $C$-modules $J_\lambda$.  Their validity is exactly
where the square-zero hypothesis is used.  The compatibility
equations $q_{\lambda\mu}w_\mu=w_\lambda$ are affine linear of
the same kind.  Every finite subsystem has a solution: choose an
index $\nu$ dominating the finitely many indices and lift
$u_\nu$ by formal smoothness, then project this lift to those
indices.  Lemma~\ref{SA23-D-1} supplies compatible lifts, hence an
$A$-algebra map $v:B\to C$ reducing to $u$; reduction is justified
by Lemma~\ref{SA23-D-2} with $D=J$.

This algebra map is continuous.  For each $\lambda$, continuity of
$u_\lambda$ gives an open ideal $V\subset B$ with
$u_\lambda(V)=0$.  Its image under $w_\lambda$ lies in
$J_\lambda$, so $w_\lambda(V^2)=0$.  The ideal $V^2$ is open by
hypothesis.  Thus every $w_\lambda$ is continuous, and so is $v$.

For general $J$, put $J_n=\overline{J^n}$ for $n\geq1$; in
particular $J_1=J$.  Continuous multiplication gives
$J_n^2\subset\overline{J^{2n}}\subset J_{n+1}$.  By
Lemma~\ref{SA23-D-2}, each $C/J_n$ is complete and pro-Artinian.
The map $C/J_{n+1}\to C/J_n$ has closed square-zero kernel
$J_n/J_{n+1}$.  Applying the square-zero case repeatedly constructs
compatible continuous lifts $u_n:B\to C/J_n$ with $u_1=u$.

We finally verify the comparison used to take their limit.  For
each $\lambda$, nilpotence of the image of $J$ implies
$J^{n(\lambda)}\subset U_\lambda$, hence
$J_{n(\lambda)}\subset U_\lambda$ since $U_\lambda$ is closed.
Thus $\bigcap_nJ_n=0$ and
\[
 C\longrightarrow\varprojlim_{n\geq1}C/J_n
\]
is injective.  To show surjectivity choose representatives $c_n$
of compatible classes.  For $m\geq n$ their difference lies in
$J_n$.  Modulo each $U_\lambda$ the sequence is therefore
eventually constant; these eventual residues form an element of
$\varprojlim C_\lambda=C$.  Its difference from $c_n$ belongs
to $J_n$ by closedness.  This proves surjectivity.  The comparison
is a homeomorphism: any open $U_\lambda$ contains some $J_n$,
and $U_\lambda/J_n$ is open in $C/J_n$.  The inverse limit of
the continuous maps $u_n$ consequently gives the required
continuous map $v:B\to C$.
\end{proof}

\begin{proposition}[Refuted proof step; no counterexample to the theorem]\label{SA23-D-4}
The proof printed in $0_{\rm IV}$ 19.3.12, pp.~85--86, wrongly
treats sets of lifts through arbitrary nilpotent ideals as affine
linear varieties.  Deligne's polynomial compactness counterexample
on cote 23, scan 060, is correct, but it does not refute
Theorem~\ref{SA23-D-3}.
\end{proposition}
\begin{proof}
For an infinite field $k$ consider variables $x$ and $y_a$ for
$a\in k$, with equations $(x-a)y_a=1$.  Any finite collection is
soluble: choose $x$ different from the finitely many specified
$a$ and set $y_a=(x-a)^{-1}$.  A simultaneous solution would
require $x\ne a$ for every $a\in k$, which is impossible.
This is a counterexample to the polynomial compactness principle
used by the printed proof.  In a lift $w^0+z$ the multiplicativity
equation contains $z(b)z(b')$ when the kernel is not square-zero,
so it need not be affine linear.  The repaired proof applies linear
compactness only to square-zero kernels, then takes the countable
tower of closed powers.  The statement therefore remains valid.
The literature searches recorded in the source log found the printed
proposition and its mirrors, but no prior source for this particular
repair; we make no priority claim.
\end{proof}

\subsubsection{D.2. The filtered and graded assertions}

\begin{proposition}[Proved: IV 19.5.5, the implication in the letter]\label{SA23-D-5}
Let $R$ be a ring, $M=F_0M\supset F_1M\supset\cdots$ a filtered
$R$-module, $f=(f_1,\ldots,f_r)$, and $I=(f)$.  Put
$G_nM=\sum_{i=0}^n I^iF_{n-i}M$ and $X=\operatorname{gr}_F M$.
If the canonical graded surjection
\[
 \psi:(X/IX)[T_1,\ldots,T_r]\longrightarrow
       \operatorname{gr}_G M,
 \qquad [m]T^\alpha\longmapsto[f^\alpha m]
\]
is an isomorphism, then $f$ is $X$-quasi-regular.  Here $[m]\in X_p$
and $T^\alpha$ has degree $|\alpha|$, so the target degree is
$p+|\alpha|$.  No separatedness assumption is needed.
This is the final printed IV 19.5.5(b)$\Rightarrow$(c), pp.~206--209;
the footnote on p.~207 explicitly credits Deligne's proof.  The
corresponding numbers in $0_{\rm IV}$, p.~94, are a different remark
and a proof of Theorem 19.5.3, not the claims criticized in this letter.
\end{proposition}
\begin{proof}
Injectivity of $\psi$ has this useful consequence: if a finite sum
\[
 \sum_{p+|\alpha|=n} f^\alpha m_{p,\alpha}\in G_{n+1}M,
 \qquad m_{p,\alpha}\in F_pM,
\]
then every $m_{p,\alpha}$ lies in $F_{p+1}M+IF_pM$.
This is simply the direct-sum decomposition of the degree-$n$
part of the source of $\psi$.

We prove by induction on $q\geq1$ that
\[
 I^qF_pM\cap F_{p+1}M
 \subset\sum_{i=1}^q I^{q-i}F_{p+i}M
 \quad\hbox{for every }p\geq0.                         \tag{D.5.1}
\]
The case $q=1$ is immediate.  For
$m\in I^{q+1}F_pM\cap F_{p+1}M$, the induction hypothesis writes
$m=\sum_{i=1}^q\sum_{|\alpha|=q-i}f^\alpha m_{i,\alpha}$,
where $m_{i,\alpha}\in F_{p+i}M$.  But
$m\in I^{q+1}F_pM\subset G_{p+q+1}M$.  The displayed consequence
of injectivity, in total degree $p+q$, gives
$m_{i,\alpha}\in F_{p+i+1}M+IF_{p+i}M$ for every $i,\alpha$.
Substitution yields (D.5.1) for $q+1$.

To prove quasi-regularity take $x_\alpha\in F_pM$, $|\alpha|=q$,
such that
$\sum f^\alpha x_\alpha\in F_{p+1}M+I^{q+1}F_pM$.
Subtracting suitable elements of $IF_pM$ from the $x_\alpha$
makes this sum belong to $F_{p+1}M$ without changing their classes
in $F_pM/(F_{p+1}M+IF_pM)$.  Formula (D.5.1) then expresses the
sum using terms $f^\beta y_{i,\beta}$ with
$y_{i,\beta}\in F_{p+i}M$ and $|\beta|=q-i$, $i>0$.
The resulting equality, in total degree $p+q$, and injectivity of
$\psi$ give $x_\alpha\in F_{p+1}M+IF_pM$.  This is exactly
injectivity of
$(X/IX)[T]_q\to I^qX/I^{q+1}X$ on each $F$-degree.  Surjectivity
is automatic from the generators of $I$, proving the assertion.
\end{proof}

\begin{example}[Refuted: the final IV 19.5.6(ii); corrected scope]\label{SA23-D-6}
The final 1967 text, p.~209, assumes only that the $f_i$ have no
homogeneous component of degree zero.  Under that wording its
quasi-regular$\Rightarrow$regular assertion is false.  For any field
$k$, take the ordinarily graded ring and module
\[
 R=M=k[x,y,z],\qquad f=((1-x)y,(1-x)z,x).
\]
Every $f_i$ has strictly positive degree components, $M$ is bounded
below, and $f$ is Koszul-regular and quasi-regular.  It is not an
ordinary regular sequence.  A valid corrected statement is obtained
by requiring the $f_i$ to be homogeneous of positive degree; more
generally the exact sufficient hypothesis is separatedness of each
successive quotient for the $I$-adic topology, as in
$0_{\rm IV}$ 15.1.9.
\end{example}
\begin{proof}
The sequence generates $(x,y,z)$, since
$y=(1-x)y+xy$ and $z=(1-x)z+xz$.  The latter ideal is generated
by the ordinary regular sequence $(x,y,z)$.  The generator-invariance
and quasi-regularity results of
Theorems~\ref{SA23-C-3} and \ref{SA23-C-2} show that $f$ is
Koszul-regular and quasi-regular.  Its first element is a
nonzerodivisor, because $R$ is a domain.  In
$R/((1-x)y)$, however, the nonzero class of $y$ is killed by
$(1-x)z$, since $y(1-x)z=z(1-x)y$.
The class of $y$ is nonzero: cancellation in the domain would
otherwise make $1-x$ a unit in $k[x,y,z]$.  Thus the second
successive injectivity condition fails.

For completeness, the separation criterion needed for the correction
follows directly from quasi-regularity.  If
$(M/IM)[T]\simeq\operatorname{gr}_I M$, then multiplication by
$f_1$ is strict and injective whenever $M$ is $I$-adically separated:
if $f_1m\in I^{n+1}M$, comparison of its initial form shows
$m\in I^nM$ by induction on its filtration order; for
$f_1m=0$ this puts $m$ in every $I^nM$.  Strictness means
$f_1M\cap I^nM=f_1 I^{n-1}M$ for $n\geq1$.
It identifies the associated graded of $M/f_1M$ with
$(M/IM)[T_2,\ldots,T_r]$.  Repeating this argument, assuming
separatedness of the successive quotients, proves regularity.
When the $f_i$ are homogeneous of positive degree, these quotients
remain graded and bounded below; $I^n$ raises the least possible
degree by at least $n$, proving the required separatedness.
No finite generation of the graded module is used.  If one's
definition of regularity requires a nonzero final quotient, it is
included as a separate hypothesis (and holds in the example).

The earlier, stronger claim that \emph{every} subquotient is separated
also fails even with homogeneous generators: for $R=k[T]$ and
$I=(T)$, the quotient $R/(T-1)$ is nonzero and equals its own
$I$-multiple.  Likewise the relative Noetherian finiteness mentioned
in Deligne's December letter is insufficient for arbitrary
$A$-subquotients: take $A=k[x]_{(x)}$, $B=A[[T]]$, $M=B$.
The $A$-submodule $A[T]\subset B$ has quotient
$A[T]/(xT-1)\simeq k(x)$, on which $x$ is invertible.  Restricting
to finite $B$-subquotients and $IB\subset\operatorname{Jac}(B)$
does give separatedness by the Noetherian Krull intersection theorem.
\end{proof}
